\documentclass[8pt, a4paper, landscape]{extarticle}
\usepackage[utf8]{inputenc}
\usepackage[russian]{babel}
\usepackage{amsmath, amssymb, amsthm}
\usepackage[margin=0.4cm]{geometry}
\usepackage{multicol}
\usepackage{enumitem}
\usepackage{sectsty}

\sectionfont{\normalsize\bfseries\scshape\vspace{-2ex}\hrule\vspace{1ex}}
\subsectionfont{\small\bfseries\vspace{-2ex}}
\setlength{\parindent}{0pt}
\setlength{\parskip}{1pt}
\setlength{\abovedisplayskip}{1pt}
\setlength{\belowdisplayskip}{1pt}
\setlist{nosep, leftmargin=1.2em}

\newcommand{\E}{\mathbb{E}}
\renewcommand{\P}{\mathbb{P}}
\DeclareMathOperator{\Var}{Var}
\DeclareMathOperator{\Cov}{Cov}

\begin{document}
\begin{multicols*}{4}

\section*{1. Асимптотика и ЛБК}
\textbf{Леммы Бореля-Кантелли (ЛБК):}
\begin{enumerate}
    \item $\sum \P(A_n) < \infty \implies \P(A_n \text{ б.ч.}) = 0$.
    \item $\sum \P(A_n) = \infty$, $A_n$ незав. $\implies \P(A_n \text{ б.ч.}) = 1$.
\end{enumerate}

\textbf{Алгоритм поиска предела $C = \varlimsup \frac{\xi_n - a_n}{b_n}$:}
\begin{enumerate}
    \item \textbf{Порог:} Вводим $\varepsilon>0$, события $A_n^\pm = \{ \eta_n > C \pm \varepsilon \}$.
    \item \textbf{Хвосты:} $\P(A_n) = \P(\xi_n > a_n + b_n(C \pm \varepsilon))$. Заменяем вероятность на асимптотику хвоста.
    \item \textbf{Ряды:} Подстановка дает обобщенный гарм. ряд $\sum \frac{(\ln n)^\beta}{n(\ln n)^C} \propto \sum \frac{1}{n(\ln n)^q}$. Сходится при $q>1$, расходится при $q \le 1$.
    \item \textbf{Синтез:} Подбираем $C$, чтобы для $A_n^+$ ряд сходился (ЛБК-1 $\implies \varlimsup \le C$), а для $A_n^-$ расходился (ЛБК-2 $\implies \varlimsup \ge C$).
\end{enumerate}

\textbf{Асимптотика хвостов $\P(\xi > x)$ при $x \to \infty$:}
\begin{itemize}
    \item $N(0,1): \sim \frac{1}{\sqrt{2\pi} x} e^{-x^2/2}$ (убывает как $x^{-1}$)
    \item $\Gamma(\alpha, \lambda): \sim \frac{\lambda^{\alpha-1}}{\Gamma(\alpha)} x^{\alpha-1} e^{-\lambda x}$
    \item $\text{Exp}(\lambda): = e^{-\lambda x}$ (точное равенство)
    \item Тяж. хвосты (Коши): $\sim C x^{-\alpha}$. Нормировка $b_n=n^{1/\alpha}$.
\end{itemize}

\textbf{Мастер-теорема для $K x^\beta e^{-\lambda x}$:}
При стандартной нормировке $a_n = \frac{1}{\lambda}\ln n, b_n = \frac{1}{\lambda}\ln\ln n$, граница сходимости (искомый предел): $\mathbf{C = \beta + 1}$.
$\text{Exp}(\lambda): \beta=0 \Rightarrow C=1$. 
$\Gamma(\alpha, \lambda): \beta=\alpha-1 \Rightarrow C=\alpha$. 
$N(0,1):$ $a_n=\sqrt{2\ln n}$, $\beta=-1 \Rightarrow C=-1/2$.

\section*{2. Сходимости, ЦПТ и Дельта-метод}
\textbf{Топология сходимостей:}
\begin{itemize}
    \item $\xi_n \xrightarrow{d} C \iff \xi_n \xrightarrow{\P} C$ (\textbf{только} если предел --- константа!).
    \item ЦПТ: $\eta_n = \sqrt{n}(\overline{X} - \mu)$ сходится по $d$, но \textbf{не имеет} предела по $\P$ (разность $\eta_n - \eta_m$ не стремится к нулю, не фундаментальна).
    \item Вектора: $\xi_n \xrightarrow{d} \xi, \eta_n \xrightarrow{d} \eta \not\Rightarrow (\xi_n, \eta_n) \xrightarrow{d} (\xi, \eta)$. Нужна независимость компонент или многомерная ЦПТ.
\end{itemize}

\textbf{Алгоритм Дельта-метода (Сходимость по $d$):}
\begin{enumerate}
    \item \textbf{База:} Найти параметры $X_1$: $\mu = \E X_1$, $\sigma^2 = \Var X_1$.
    \item \textbf{ЦПТ:} Записать базу: $\sqrt{n}(\overline{X} - \mu) \xrightarrow{d} N(0, \sigma^2)$.
    \item \textbf{Функция:} Выделить $g(t)$. Проверить, что вычитаемая в условии константа строго равна $C = g(\mu)$.
    \item \textbf{Дисперсия:} Найти $g'(\mu)$. Асимптотическая дисперсия $\sigma_{as}^2 = [g'(\mu)]^2 \sigma^2$. Ответ: $N(0, \sigma_{as}^2)$.
\end{enumerate}

\textbf{ВЫРОЖДЕННЫЙ СЛУЧАЙ (2-й порядок, Теорема о $\chi^2$):}
Если $g'(\mu)=0$, но $g''(\mu)\ne 0$, классический метод дает $\sigma_{as}^2=0$ ($\xrightarrow{\P} 0$). Домножаем ряд Тейлора на $n$ (а не на $\sqrt{n}$):
$ n(g(\overline{X}) - g(\mu)) \xrightarrow{d} \frac{\sigma^2 g''(\mu)}{2} \chi^2_1 $
\textit{Предел — гамма-распределение (пропорционально $\chi^2$), а не норм!}

\textbf{Многомерный Дельта-метод:} $\vec{W}_i=(Y_i, Z_i)^T$ --- выборка, $\E[\vec{W}_1]=\vec{\mu}$, $\Sigma = \Cov(\vec{W}_1)$. Формула: $\sigma_{as}^2 = \nabla g(\vec{\mu})^T \Sigma \nabla g(\vec{\mu})$.
\textit{Синтез для $g(y,z) = z/y^k$:} Градиент $\nabla g = \left(-\frac{k\mu_z}{\mu_y^{k+1}}, \frac{1}{\mu_y^k}\right)^T$.
$ \sigma^2_{as} = \frac{k^2 \mu_z^2}{\mu_y^{2k+2}} \Var(Y) - \frac{2k \mu_z}{\mu_y^{2k+1}} \Cov(Y, Z) + \frac{1}{\mu_y^{2k}} \Var(Z) $

\section*{3. Гауссовские векторы}
\textbf{Алгоритм вычисления нелинейного $\E[\cos(YZ)|X]$:}
\begin{enumerate}
    \item \textbf{Проекции:} Представить $Y = aX+\tilde{Y}$, $Z=bX+\tilde{Z}$.
    $a \!=\! \frac{\Cov(Y,X)}{\Var(X)}$, $\Var(\tilde{Y}) \!=\! \Var(Y) - \frac{\Cov^2(X,Y)}{\Var(X)}$.
    Вычислить $\Cov(\tilde{Y},\tilde{Z}) \!=\! \Cov(Y,Z) - ab\Var(X)$. Если 0 $\implies \tilde{Y} \perp\!\!\perp \tilde{Z}$. Зафиксируем $X=x \implies Y_x, Z_x$ независимы.
    \item \textbf{Комплексификация:} Интеграл косинуса не берется. Формула Эйлера: $\cos \alpha = \text{Re}(e^{i\alpha})$. Ищем $\text{Re}(\E[e^{iY_x Z_x}])$.
    \item \textbf{Итерированное $\E$:} По свойству башни обусловим по $Z_x$. Внутреннее $\E$ --- это ХФ нормальной $Y_x$:
    $\E[e^{i Z_x Y_x}|Z_x] = \exp(i\mu_{Y_x} Z_x - \frac{1}{2}\sigma_{Y_x}^2 Z_x^2)$.
    \item \textbf{Интеграл Гаусса:} Внешнее $\E$ по $Z_x \sim N(\mu_{Z_x}, \sigma_{Z_x}^2)$ (с константами $c=i\mu_{Y_x}, d=-\frac{1}{2}\sigma_{Y_x}^2$):
    $ \E[e^{c\eta + d\eta^2}] = \frac{1}{\sqrt{1-2d\sigma^2}}\exp\left( \frac{c^2\sigma^2+2c\mu+2d\mu^2}{2(1-2d\sigma^2)} \right) $
\end{enumerate}

\textbf{Формула Вика:} Для $\vec{X} \sim N(\vec{0}, \Sigma)$: $\E[\text{нечетн.}]=0$.
$\E[X_1 X_2 X_3 X_4] = \E[X_1 X_2]\E[X_3 X_4] + \E[X_1 X_3]\E[X_2 X_4] + \E[X_1 X_4]\E[X_2 X_3]$.

\section*{4. Смешанные распределения}
$X$ имеет скачки $D=\{x_k\}$ и непрерывные участки. Обобщенная плотность: $p_X^*(x) = F'_{ac}(x) + \sum p_k \delta(x-x_k)$.
\textbf{Алгоритм $\E[g(X,Y)]$ (Смешанная $X \perp\!\!\perp$ Непрерывн. $Y$):}
\begin{enumerate}
    \item \textbf{Внутреннее $\E$ (Тонелли):} Сначала берем $\E$ по непрерывной $Y$! Фиксируем $X=x$. $\psi(x) = \E[g(x,Y)] = \int g(x,y)p_Y(y)dy$.
    \item \textbf{Атомы $X$:} Точки разрыва $x_k$. Веса: $p_k = F_X(x_k) - F_X(x_k-0)$.
    \item \textbf{Непрерывная $X$:} Вне разрывов $p_{ac}(x) = F'_X(x)$. \textit{Проверка нормировки:} строго $\int p_{ac}dx + \sum p_k = 1$.
    \item \textbf{Синтез:} $\E[g(X,Y)] = \int_{\mathbb{R}\setminus D} \psi(x) F'_{ac}(x) dx + \sum_{k} \psi(x_k) p_k$.
\end{enumerate}

\textbf{БЫСТРАЯ ФОРМУЛА для показательной функции $\E[Y^X]$:}
Если $Y \sim U(a,b)$, внутреннее ожидание интегрируется мгновенно:
$ \psi(x) = \E[Y^x] = \int_a^b y^x \frac{1}{b-a} dy = \frac{b^{x+1}-a^{x+1}}{(b-a)(x+1)} $

\section*{5. Условные плотности и Якобиан}
\textbf{Алгоритм метода Якобиана (поиск $p_{X|Z}$ где $Z=f(X,Y)$):}
\begin{enumerate}
    \item \textbf{Совместная:} Если $X \perp\!\!\perp Y \implies p_{X,Y}(x,y) = p_X(x)p_Y(y)$.
    \item \textbf{Замена:} Вводим $U=X, Y=h(X,Z)$. Вычисляем якобиан $J = \partial h(x,z) / \partial z$. Плотность: $p_{X,Z} = p_{X,Y}(x, h(x,z)) \cdot |J|$.
    \item \textbf{ДИНАМИЧЕСКИЕ ГРАНИЦЫ (Критично!):} Подставляем $y=h(x,z)$ в индикаторы исходной плотности $I_{\{y \in D_y\}}$. Разрешаем систему относительно $x$. \textit{Пример:} $x \in (0,1), xz \in (0,1) \implies 0 < x < \min(1, 1/z)$.
    \item \textbf{Маргинализация:} Интегрируем $p_{X,Z}$ по $x$. Из-за $\min$ интеграл $p_Z(z)$ распадается на ветви (напр. $z \le 1$ и $z > 1$).
    \item \textbf{Усл. плотность и ожидание:} $p_{X|Z} = p_{X,Z}/p_Z$. Вычисляем $\E[g|Z=z] = \int g(\dots) p_{X|Z}(x|z)dx$.
\end{enumerate}

\textbf{Трюк измеримости (Вынесение множителя):}
Относительно условия $Z$, сама $Z$ — константа.
$\E[XY | Z=Y/X] = \E[X(XZ)|Z] = Z \cdot \E[X^2|Z]$. (Ищем только $\E[X^2|Z]$!)

\section*{Комбинаторика и Симметрия}
\textbf{Комбинаторная симметрия:} Если $X_1 \dots X_n$ н.о.р., а функция условия $S = \sum X_i$ симметрична, то все $\E[X_k|S]$ равны.
$\sum \E[X_k|S] = \E[S|S] = S \implies \E[X_k|S] = S/n$. Без интегралов!

\textbf{Немонотонные условия:} Если биекции нет ($Z=X^2$).
$X \sim U(-a,a)$, плотность четная. Усл. распр. дискретно: $\P(X=\sqrt{z}|z) = \P(X=-\sqrt{z}|z) = 1/2 \implies \E[X|X^2=z] = 0$.

\section*{Характеристические функции}
$\varphi(t) = \E[e^{it\xi}]$. Свойства: $\varphi(0)=1$, $|\varphi(t)| \le 1$, $\varphi_{aX+b}(t) = e^{itb}\varphi_X(at)$.
$\xi \perp\!\!\perp \eta \implies \varphi_{\xi+\eta} = \varphi_\xi \varphi_\eta$. Моменты: $\E[X^k] = \varphi^{(k)}(0) / i^k$.
\textbf{Решетчатость:} $|\varphi(t_0)|=1$ для $t_0 \ne 0 \iff$ распр. дискретное.
\textbf{Таблица ХФ:}
\begin{itemize}
    \item $N(\mu,\sigma^2): \exp(i\mu t - \frac{1}{2}\sigma^2 t^2)$ \quad $\text{Po}(\lambda): \exp(\lambda(e^{it}-1))$
    \item $\text{Bin}(n,p): (1-p+pe^{it})^n$ \quad $\text{Exp}(\lambda): \lambda / (\lambda-it)$
    \item $U(a,b): (e^{itb}-e^{ita}) / (it(b-a))$
    \item Лаплас $L(0,b): 1 / (1+b^2 t^2)$ \quad Коши $K(a): e^{-a|t|}$
\end{itemize}

\section*{Случайные суммы и Неравенства}
\textbf{Чернов:} $\P(X \ge a) \le \inf_{t>0} e^{-ta}\E[e^{tX}]$.
\textbf{Случайные суммы:} $S_\nu = \sum_{i=1}^\nu \xi_i$ ($\nu \perp\!\!\perp \xi_i$).
\textbf{Вальд:} $\E[S_\nu] \!=\! \E\nu\E\xi_1$. $\Var(S_\nu) \!=\! \E\nu\Var\xi_1 \!+\! \Var\nu(\E\xi_1)^2$.
Если $\nu \sim \text{Po}(\lambda)$, $\xi_i \sim \text{Be}(p)$, то успехи $S_\nu \sim \text{Po}(\lambda p)$ и неудачи $\nu-S_\nu \sim \text{Po}(\lambda(1-p))$ --- \textbf{независимы}!

\section*{Базовые моменты и Быстрый счет}
Фундаментальное тождество: $\E[X^2] = \Var X + (\E X)^2$. 
\textbf{Симметрия модулей:} Для $X \sim N(0, \sigma^2)$ или Лапласа $\E[X^{\text{нечет}}]=0$. Ковариация: $\Cov(|X|, X^2) = \E[|X|^3] - \E|X|\E[X^2]$. \textbf{Не нужно раскрывать модуль!} Вычисляется в 1 строку.
\textbf{Абсолютные моменты $N(0, \sigma^2)$:}
$\E|X|^p = \sigma^p 2^{p/2} \Gamma((p+1)/2)/\sqrt{\pi}$.
$\E|X| = \sigma \sqrt{2/\pi}, \quad \E[X^2]=\sigma^2, \quad \mathbf{\E|X|^3 = 2\sigma^3 \sqrt{2/\pi}}$.

\textbf{Сводка Распределений ($\E X, \Var X, \E[X^2]$):}
\begin{itemize}
    \item \textbf{Be($p$):} $p, \ p(1-p), \ p$
    \item \textbf{Bin($n,p$):} $np, \ npq, \ npq+n^2p^2$
    \item \textbf{Po($\lambda$):} $\lambda, \ \lambda, \ \lambda(\lambda+1)$
    \item \textbf{Geom($p$) [от 1]:} $1/p, \ (1-p)/p^2, \ (2-p)/p^2$
    \item \textbf{$U(a,b)$:} $\frac{a+b}{2}, \ \frac{(b-a)^2}{12}, \ \frac{a^2+ab+b^2}{3}$
    \item \textbf{$N(\mu,\sigma^2)$:} $\mu, \ \sigma^2, \ \mu^2+\sigma^2$ (для $N(0,\sigma^2)$: $\E[X^4]=3\sigma^4$)
    \item \textbf{Exp($\lambda$):} $1/\lambda, \ 1/\lambda^2, \ 2/\lambda^2$ (пл. $\lambda e^{-\lambda x}$)
    \item \textbf{$\Gamma(\alpha,\lambda)$:} $\alpha/\lambda, \ \alpha/\lambda^2, \ \alpha(\alpha+1)/\lambda^2$
    \item \textbf{Лаплас(0,$b$):} $0, \ 2b^2, \ 2b^2$ \quad (пл. $\frac{1}{2b}e^{-|x|/b}$)
    \item \textbf{Коши $K(x_0, \gamma)$:} $\E X, \Var X$ \textbf{не существуют}! ЗБЧ и ЦПТ неприменимы. Выборочное среднее $\overline{X} \sim K(x_0, \gamma)$.
\end{itemize}

\textbf{Sanity Checks (Самопроверка):}
\begin{enumerate}
    \item \textbf{Размерность:} Если $[X]\!=\!\text{м}$, то $[\Var X]\!=\!\text{м}^2$. Для $g(X)\!=\!1/X^2$, дисперсия $\sigma_{as}^2 \sim [g]^2 \!=\! \text{м}^{-4}$. Нарушение = ошибка дифференцирования.
    \item \textbf{Квадратичная форма:} Если $\sigma_{as}^2 = \nabla g^T \Sigma \nabla g < 0$ --- ошибка в знаках ковариации или Якобиане.
    \item \textbf{Тест константы:} В условном ожидании $\E[g(X,Y)|Z=z] = \varphi(z)$. Если $m \le g \le M$, то и $m \le \varphi(z) \le M$.
\end{enumerate}

\end{multicols*}
\end{document}
